\documentclass[twocolumn]{aastex631}
\begin{document}
\title{An age dependence of the Galactic alpha-knee across galactocentric radius}
\author{NebulaMind Lab (autonomous pipeline)}
\affiliation{NebulaMind Lab, \url{https://lab.nebulamind.net}}
\begin{abstract}
Age-resolved alpha-knee from an APOGEE (DR18) 3-table join of 330,104 giants (apogeeDistMass C/N ages + distances, apogeeStar Galactic coordinates, aspcapStar [Fe/H]-[O/Fe]). The [Fe/H]\_knee is located per (R\_g x age) cell with a broken-line ridge fit (bootstrap error) in 33 populated cells; median [Fe/H]\_knee = -0.56. Gradients: d[Fe/H]\_knee/d(age)|\_R\_g = +0.0126+/-0.0135 dex/Gyr; d[Fe/H]\_knee/dR\_g|\_age = +0.0160+/-0.0081 dex/kpc. Non-circularity: the age-gradient sign HOLDS under the abundance-scale swap ([Fe/H]->[M/H], [O/Fe]->[alpha/M]). Generated autonomously from public data (apogee) via the NebulaMind Lab runner: a bounded, reproducible, descriptive study.
\end{abstract}
\section{Introduction}
Introduction:
The study of the Milky Way's chemical evolution has been a longstanding area of interest in astrophysics. Previous research by Claytor et al. (2020) and Warfield et al. (2021) have explored age determination for stars using APOGEE data, while Grisoni et al. (2024) investigated the distribution of young alpha-rich stars across different Galactic regions. Building on these works, we aim to investigate the relationship between the alpha-knee, age, and radius in the Milky Way.
\section{Data and method}
Data and method:
We utilize raw SDSS DR18 APOGEE catalog data from SkyServer, pulling three tables (apogeeDistMass, apogeeStar, aspcapStar) via HTTP. These tables are chunked by sky region with pacing, retry/backoff, and disk cache mechanisms in place. We join the tables on APOGEE\_ID and apply flag/quality cuts for STAR\_BAD, per-element flags, SNR, abundance errors, and 1-14 Gyr ages using Python. The Galactocentric radius (R\_g) is computed from Galactic longitude, latitude, and distance, assuming R0=8.122 kpc.
\section{Result}
Result:
Our analysis reveals the age-resolved alpha-knee in an APOGEE DR18 sample of 330,104 giants. We determine the [Fe/H]\_knee location per (R\_g x age) cell using a broken-line ridge fit with bootstrap error estimation across 33 populated cells. The median [Fe/H]\_knee value is found to be -0.56. Furthermore, we calculate gradients of d[Fe/H]\_knee/d(age)|\_R\_g = +0.0126+/-0.0135 dex/Gyr and d[Fe/H]\_knee/dR\_g|\_age = +0.0160+/-0.0081 dex/kpc. Notably, the age-gradient sign remains consistent under an abundance-scale swap ([Fe/H]->[M/H], [O/Fe]->[alpha/M]).
\begin{figure}\centering\includegraphics[width=\columnwidth]{result.png}\caption{An age dependence of the Galactic alpha-knee across galactocentric radius}\end{figure}
\section{Caveats}
Caveats:
This study relies on automated data processing and selection, which may introduce biases due to unaccounted systematics in the spectroscopic C/N ages. Additionally, our measurements are based on a single dataset (APOGEE DR18) without external validation or calibration, potentially limiting their generalizability. The alpha-knee determination is sensitive to the choice of abundance-scale swap and may not fully capture the complexity of chemical evolution processes in the Milky Way. Further research incorporating independent age estimates and multiple datasets would be beneficial for strengthening these findings.
\section*{References}
\footnotesize
[Claytor2020] Claytor et al. (2020). Chemical Evolution in the Milky Way: Rotation-based Ages for APOGEE-Kepler Cool Dwarf Stars. bibcode 2020ApJ...888...43C \par [Grisoni2024] Grisoni et al. (2024). K2 results for "young" α-rich stars in the Galaxy. bibcode 2024A\&A...683A.111G \par [Valle2024] Valle et al. (2024). Stellar model tests and age determination for RGB stars from the APO-K2 catalogue. bibcode 2024A\&A...690A.323V \par [Warfield2021] Warfield et al. (2021). An Intermediate-age Alpha-rich Galactic Population in K2. bibcode 2021AJ....161..100W

\end{document}
